\section{Affine Connections}

\subsection{Levi-Civita Connections}
When it comes to connections, most of us probably think first of 
the Levi-Civita connection in Riemannian geometry. 
So let us first review its most basic definition, 
and then try to rewrite everything we are familiar with in the language of differential forms.
\subsubsection{Covariant Derivative}
\begin{definition}
    A \textbf{connection} $\nabla$ on a mannifold $M$ is a map
    \begin{align*}
        \nabla: C^\infty(M,TM) \times C^\infty(M,TM) &\to C^\infty(M,TM) \\
        (Y,X) &\mapsto \nabla_Y X
    \end{align*}
    satisfying the following properties:
    % \begin{align*}
    %     \nabla_{f_1Y_1 + f_2Y_2} X =& f_1 \nabla_{Y_1} X + f_2 \nabla_{Y_2} X, \\
    %     \nabla_{Y} (\lambda_1X_1 + \lambda_2 X_2) =& \lambda_1 \nabla_{Y} X_1 + \lambda_2 \nabla_{Y} X_2, \\
    %     \nabla_{Y} fX =& \md f(Y)\cdot X + f \nabla_{Y} X. 
    %  \end{align*}
    \begin{flalign*}
        & \text{(C1)} && \nabla_{f_1Y_1 + f_2Y_2} X = f_1 \nabla_{Y_1} X + f_2 \nabla_{Y_2} X, &\\
        & \text{(C2)} && \nabla_{Y} (\lambda_1X_1 + \lambda_2 X_2) = \lambda_1 \nabla_{Y} X_1 + \lambda_2 \nabla_{Y} X_2, &\\
        & \text{(C3)} && \nabla_{Y} fX = \md f(Y)\cdot X + f \nabla_{Y} X. &
    \end{flalign*}
    where $X,X_1,X_2,Y,Y_1,Y_2 \in \Gamma(TM)$, $f,f_1,f_2 \in C^\infty(M)$, 
    and $\lambda_1,\lambda_2 \in \mathbb{R}$.
\end{definition}




\begin{definition}
    The \textbf{torsion}
    of a connection $\nabla$ is defined as the tensor
    \begin{equation*}
        T(X,Y) := \nabla_X Y - \nabla_Y X - [X,Y].
    \end{equation*}
    The \textbf{curvature} of a connection $\nabla$ is defined as the tensor
    \begin{equation*}
        R(X,Y)Z := \nabla_X \nabla_Y Z - \nabla_Y \nabla_X Z - \nabla_{[X,Y]} Z.
    \end{equation*}
\end{definition}

\begin{remark}
    The torsion $T$ can only be defined for connections on the tangent bundle, 
    since it involves both $\nabla_X Y$ and $\nabla_Y X$.
\end{remark}

\begin{definition}
    A connection $\nabla$ of a Riemannian manifold $(M,g)$ is called a \textbf{Levi-Civita connection} 
    if it is torsion-free and compatible with the metric $g$, i.e. 
    \begin{flalign*}
        & \text{(LC1)} && \nabla_X Y - \nabla_Y X = [X,Y], &\\
        & \text{(LC2)} && X(g(Y,Z)) = g(\nabla_X Y, Z) + g(Y, \nabla_X Z).
    \end{flalign*}
\end{definition}

\subsubsection*{Exterior convariant derivative}
A connection $\nabla$ can also be viewed as a map
\begin{align*}
    \nabla: C^\infty(M,TM) &\to C^\infty(M,T^*M \otimes TM) \\
    X &\mapsto \nabla X
\end{align*}
where $\nabla X$ is a $1$-form taken value in $TM$\footnote{One can simply view $\nabla X$ as a $(1,1)$-tensor field.} 
and $(\nabla X)(Y) := \nabla_Y X$.
$\nabla$ is not a $C^\infty(M)$-linear map, instead, 
\begin{equation*}
    \nabla(fX) = \md f \otimes X + f \nabla X.
\end{equation*}
We can extend $\nabla$ to the exterior covariant derivative 
\begin{align*}
    \nabla: \Omega^*(M,TM) &\to \Omega^{*+1}(M,TM) \footnotemark \\
    \eta &\mapsto \nabla \eta,
\end{align*}
\footnotetext{$\Omega^r(M,TM):=C^\infty(M, \wedge^r T^*M \otimes TM)$.}
satisfying the Leibniz rule:
\begin{align*}
    \nabla(\alpha\wedge\eta) = \md\alpha\wedge\eta + (-1)^{q}\alpha\wedge \nabla\eta,
    \quad\forall \alpha\in \Omega^q(M),\eta\in \Omega^*(M,TM),
\end{align*}
For example, let $\eta=\alpha\otimes X$ be a section of $\Omega^r(M,TM)$, then
\begin{equation*}
    \nabla \eta = \nabla(\alpha\otimes X) = \md\alpha \otimes X + (-1)^r \alpha\wedge \nabla X.
\end{equation*}

Under the above setup, we can define $\nabla^2=\nabla\circ\nabla,\nabla^3,\dots,\nabla^k$, 
and we can prove the following property:
\begin{proposition}
    $\nabla^2: \Omega^*(M,TM)\to\Omega^{*+2}(M,TM)$ is an $\Omega^*(M)$-linear map, i.e. 
    for any $\alpha\in \Omega^*(M),\eta\in \Omega^*(M,TM)$, 
    \begin{equation*}
        \nabla^2(\alpha\wedge\eta) = \alpha\wedge\nabla^2 \eta.
    \end{equation*}
    Thus $R:=\nabla^2$ defines a section of $\Omega^2(M,\mathrm{End}(TM))$, 
    which is called the \textbf{curvature} of the connection $\nabla$.
\end{proposition}
\begin{proof}
    \begin{align*}
        \nabla^2\big(\alpha\wedge\eta) =& \nabla(\md\alpha\wedge\eta + (-1)^{q}\alpha\wedge \md^\nabla\eta\big) \\
        =& \md^2\alpha\wedge\eta + (-1)^{q+1}\md\alpha\wedge \md^\nabla\eta + (-1)^q \md\alpha\wedge \nabla\eta + \alpha\wedge\nabla^2 \eta \\
        =& \alpha\wedge\nabla^2 \eta. \qedhere
    \end{align*}
\end{proof}
Unlike the usual exterior derivative $\md$, we have $\nabla^2\neq 0$ in general.
The following proposition explains why $\nabla^2$ is called the curvature of $\nabla$:
\begin{proposition}\label{prop:curvature}
    For $X,Y,Z\in C^\infty(M,TM)$, we have
    \begin{equation*}
        (\nabla^2 Z)(X,Y) = \nabla_X \nabla_Y Z - \nabla_Y \nabla_X Z - \nabla_{[X,Y]} Z.
    \end{equation*}
\end{proposition}
\begin{proof}
    We will prove it after introducing the connection 1-forms.
\end{proof}

\subsubsection{Connection 1-form}
Locally, we always have a frame. Let $\{e_i\}$ be a local frame of $TM$ on an open set $U\subset M$, 
and $\{\theta^i\}$ be the dual frame.
then we can write 
\begin{equation*}
    \nabla_{e_i}e_j = \Gamma_{ij}^k e_k,
\end{equation*}
where $\Gamma_{ij}^k$ are the Christoffel symbols of the connection $\nabla$ w.r.t. $\{e_i\}$. 
Define 
\begin{equation*}
    \omega_j^k := \Gamma_{ij}^k \theta^i,
\end{equation*}
we call $\omega_j^k$ the \textbf{connection 1-form} of $\nabla$ w.r.t. $\{e_i\}$. 
So
\begin{equation*}
    \nabla e_j = \omega_j^k \otimes e_k.
\end{equation*}
Let the matrix of 1-forms $\omega = (\omega_j^k)$\footnote{The upper index $k$ is the row index, 
and the lower index $j$ is the column index.}, 
then we can write the above equation in a more compact way:
\begin{equation*}
    \nabla (\ft{e}{1}{n}) = (\ft{e}{1}{n})\,\omega = 
    (\ft{e}{1}{n})
    \begin{pmatrix}
        \omega_1^1 & \cdots & \omega_1^n \\
        \vdots &  & \vdots \\
        \omega_n^1 & \cdots & \omega_n^n
    \end{pmatrix}
    \footnotemark
\end{equation*}
\footnotetext{This is an awkward notation. 
Since we arrange the frame as row vectors so matirx $\omega$ should multiply on the right. 
For a vector field $X$ and a differential forms $\alpha\in\Omega^q(M)$, 
we have $X\otimes\alpha = \alpha \otimes X$. 
For a general $\eta\in\Omega^p(M,TM)$, we have $\eta\otimes\alpha = (-1)^{pq} \alpha \otimes \eta$.}

For any vector field $X = X^i e_i$, we have 
\begin{align*}
    \nabla X =& \nabla\left((\ft{e}{1}{n})
    \begin{pmatrix}
        X^1 \\ \vdots \\ X^n
    \end{pmatrix}
    \right) \\ 
    =& \Big(\nabla(\ft{e}{1}{n})\Big)
    \begin{pmatrix}
        X^1 \\ \vdots \\ X^n
    \end{pmatrix} + (\ft{e}{1}{n}) 
    \begin{pmatrix}
        \md X^1 \\ \vdots \\ \md X^n
    \end{pmatrix} \\
    =& (\ft{e}{1}{n}) 
    \left( \md 
    \begin{pmatrix}
        X^1 \\ \vdots \\ X^n 
    \end{pmatrix} + 
    \omega 
    \begin{pmatrix}
        X^1 \\ \vdots \\ X^n
    \end{pmatrix}
    \right).
\end{align*}
Thus under a local frame, $\nabla$ can be expressed as 
\begin{equation*}
    \nabla = \md + \omega.
\end{equation*}

We can compute the curvature in terms of the connection 1-form $\omega$:
\begin{align*}
    \nabla^2 (\ft{e}{1}{n}) =& \nabla\big(\nabla(\ft{e}{1}{n})\big) 
    = \nabla\big((\ft{e}{1}{n})\,\omega\big) \\
    =& \big(\nabla(\ft{e}{1}{n})\big)\wedge \omega + (\ft{e}{1}{n})\, \md\omega \\
    =& (\ft{e}{1}{n})\,\omega\wedge \omega + (\ft{e}{1}{n})\, \md\omega \\
    =& (\ft{e}{1}{n})\,\big(\md\omega + \omega\wedge \omega\big).
\end{align*}
Define 
\begin{equation}\label{eq:curvature-form}
    \Omega := \md\omega + \omega\wedge \omega,
\end{equation}
then $\Omega$ is a matrix of 2-forms, and we call 
$\Omega$ the \textbf{curvature 2-form} of the connection $\nabla$ w.r.t. $\{e_i\}$. 
\begin{proposition}
    Let $R(e_i,e_j)e_k = R_{ijk}^l e_l$, then
    \begin{equation*}
        \Omega_k^l(e_i,e_j) = R_{ijk}^l.
    \end{equation*}
\end{proposition}
\begin{proof}
    Since $\Omega_k^l = \md\omega_k^l + \omega_m^l \wedge \omega_k^m$, 
    \begin{align*}
        &(\md\omega^l_k)(e_i,e_j)e_l = e_i(\omega_k^l(e_j))e_l - e_j(\omega_k^l(e_i))e_l - \omega_k^l([e_i,e_j])e_l \\
        =& \nabla_{e_i}(\omega_k^l(e_j)e_l) - \omega_k^l(e_j)\nabla_{e_i} e_l - \nabla_{e_j}(\omega_k^l(e_i)e_l) + \omega_k^l(e_i)\nabla_{e_j} e_l - \omega_k^l([e_i,e_j])e_l \\
        =& \nabla_{e_i} \nabla_{e_j} e_k - \omega_k^l(e_j)\nabla_{e_i} e_l - \nabla_{e_j} \nabla_{e_i} e_k + \omega_k^l(e_i)\nabla_{e_j} e_l - \omega_k^l([e_i,e_j])e_l,
    \end{align*}
    and 
    \begin{align*}
        (\omega_m^l \wedge \omega_k^m)(e_i,e_j)e_l =& \omega_m^l(e_i)\omega_k^m(e_j)e_l - \omega_m^l(e_j)\omega_k^m(e_i)e_l \\
        =& \omega^m_k(e_j)\nabla_{e_i} e_m - \omega^m_k(e_i)\nabla_{e_j} e_m.
    \end{align*}
    Thus 
    \begin{equation*}
        \Omega_k^l(e_i,e_j)e_l 
        = \nabla_{e_i}\nabla_{e_j}e_k - \nabla_{e_j}\nabla_{e_i}e_k 
        - \nabla_{[e_i,e_j]}e_k = R(e_i,e_j)e_k. \qedhere
    \end{equation*}
\end{proof}
\begin{corollary}
    This proposition immediately implies Prop \ref{prop:curvature} since
    \begin{equation*}
        (\nabla^2 e_k)(e_i,e_j) = \Omega_k^l(e_i,e_j)e_l = R(e_i,e_j)e_k. 
    \end{equation*}
\end{corollary}

% The connection 1-form $\omega$ is not a tensor, since it depends on the choice of the local frame.
Let $\{\tilde{e_j}\}$ be another local frame, and the transition matrix 
from $\{e_i\}$ to $\{\tilde{e_j}\}$ is $a=(a_i^j)$, i.e. 
\begin{equation*}
    (\ft{\tilde{e}}{1}{n}) = (\ft{e}{1}{n})\, a, 
\end{equation*}
then 
\begin{align*}
    (\ft{\tilde{e}}{1}{n})\tilde{\omega} =& \nabla(\ft{\tilde{e}}{1}{n}) 
    = \nabla\big((\ft{e}{1}{n})\, a\big) \\
    =& \big(\nabla(\ft{e}{1}{n})\big) a + (\ft{e}{1}{n}) \md a \\
    =& (\ft{e}{1}{n}) \omega a + (\ft{e}{1}{n}) \md a \\
    =& (\ft{\tilde{e}}{1}{n}) (a^{-1} \omega a + \md a).
\end{align*}
Thus 
\begin{equation}\label{eq:trans-formula-for-con-form}
    \tilde{\omega} = a^{-1} \md a + a^{-1} \omega a.
\end{equation}
Applying the transformation formula (\ref{eq:trans-formula-for-con-form}) to (\ref{eq:curvature-form}), we get 
\begin{align*}
    \tilde{\Omega} =& \md\tilde{\omega} + \tilde{\omega}\wedge \tilde{\omega} \\
    =& \md(a^{-1} \md a + a^{-1} \omega a) + (a^{-1} \md a + a^{-1} \omega a)\wedge (a^{-1} \md a + a^{-1} \omega a) \\
    =& \md a^{-1} \wedge \md a + \md a^{-1} \wedge \omega a + a^{-1} \md\omega\, a - a^{-1} \omega \wedge \md a \\
    &+ a^{-1} \md a \wedge a^{-1} \md a + a^{-1} \md a \wedge a^{-1} \omega a 
    + a^{-1} \omega \wedge \md a + a^{-1} \omega \wedge \omega a, \\
\end{align*}
since
\begin{equation*}
    a^{-1}a = I \implies \md a^{-1} a + a^{-1} \md a = 0 \implies \md a^{-1} = -a^{-1} \md a\, a^{-1},
\end{equation*}
we have
\begin{align*}
    \tilde{\Omega} =& -a^{-1} \md a\, a^{-1} \wedge \md a - a^{-1} \md a\, a^{-1} \wedge \omega a + a^{-1} \md\omega\, a - a^{-1} \omega \wedge \md a \\
    &+ a^{-1} \md a \wedge a^{-1} \md a + a^{-1} \md a \wedge a^{-1} \omega a + a^{-1} \omega \wedge \md a + a^{-1} \omega \wedge \omega a \\
    =& a^{-1} (\md\omega + \omega\wedge\omega) a = a^{-1} \Omega a.
\end{align*}
Thus transformation formula for curvature between $\{e_i\}$ and $\{\tilde{e}_j\}$ is given by
\begin{equation}\label{eq:trans-formula-for-curvature}
    \tilde{\Omega} = a^{-1} \Omega a.
\end{equation}

The above formula shows that although $\Omega$ isn't a globally defined form, two different curvature forms 
only differ by a conjugation. So if we apply a quantity invariant under conjugation, we obtain a globally defined form. 
This is the idea of Chern-Weil theory, which we will discuss in later section.

% For a Riemannian manifold $(M,g)$, let $\nabla$ be the Levi-Civita connection. We can take $\{e_i\}$ to be a local orthonormal frame, 
% and $\{\omega^i\}$ be the corresponding dual frame. Then
% \begin{align*}
%     &\langle \nabla_{e_k} e_i, e_j \rangle + \langle e_i, \nabla_{e_k} e_j \rangle = e_k\langle e_i, e_j \rangle = 0 \\
%     \implies &\langle \omega^l_i(e_k)e_l, e_j \rangle + \langle e_i, \omega^l_j(e_k)e_l \rangle = 0 \\
%     \implies &\omega^j_i(e_k) + \omega^i_j(e_k) = 0.
% \end{align*}
% So the connection 1-form $\omega$ w.r.t. an orthonormal frame is skew-symmetric, i.e. $\omega \in \Omega^1(M,\mathfrak{so}(n))$. 
% Since $\Omega=\md\omega+\omega\wedge\omega$, the curvature 2-form $\Omega$ is also skew-symmetric, 
% i.e. $\Omega \in \Omega^2(M,\mathfrak{so}(n))$.


\subsection{General Affine Connections}
Let $\pi:E\to M$ be a rank $m$ vector bundle, we can also define affine connections on $E$ 
in a similar way as we did for the tangent bundle. 
\subsubsection{Covariant Derivative}
The content of this section is almost a verbatim copy of the previous content. 
\begin{definition}\label{def:con-on-E}
    An (affine) connection $\nabla$ on a vector bundle $E$ is a map
    \begin{align*}
        \nabla: C^\infty(M,TM) \times C^\infty(M,E) &\to C^\infty(M,E) \\
        (X,s) &\mapsto \nabla_X s
    \end{align*}
    satisfying 
    \begin{flalign*}
        & \text{(C1)} && \nabla_{f_1Y_1 + f_2Y_2} X = f_1 \nabla_{Y_1} X + f_2 \nabla_{Y_2} X, &\\
        & \text{(C2)} && \nabla_{Y} (\lambda_1X_1 + \lambda_2 X_2) = \lambda_1 \nabla_{Y} X_1 + \lambda_2 \nabla_{Y} X_2, &\\
        & \text{(C3)} && \nabla_{Y} fX = \md f(Y)\cdot X + f \nabla_{Y} X. &
    \end{flalign*}
\end{definition}

Let $(\nabla s)(X):= \nabla_X s$, we can extend $\nabla$ to the exterior covariant derivative
\begin{align*}
    \nabla: \Omega^*(M,E) &\to \Omega^{*+1}(M,E) \\
    \eta &\mapsto \nabla \eta,
\end{align*}
satisfying the Leibniz rule: for any $\alpha\in \Omega^q(M),\eta\in \Omega^*(M,E)$,
\begin{align*}
    \nabla(\alpha\wedge\eta) = \md\alpha\wedge\eta + (-1)^{q}\alpha\wedge \nabla\eta.
\end{align*}

\begin{remark}
    The reason why we call $\nabla$ an \textbf{affine} connection is that given two connections $\nabla_1$ and $\nabla_2$ on $E$, 
    $(1-t)\nabla_1 + t\nabla_2$ is also a connection on $E$ for any $t\in \mathbb{R}$.
\end{remark}

\begin{remark}
    % For a tivial bundle $M\times \mathbb{R}^m$, we can easily define a connection $\nabla=\md$, 
    % \begin{equation*}
    %     \nabla_X (\ft{s}{1}{m})^T = (\ft{Xs}{1}{m})^T.
    % \end{equation*}
    % Thus there always exists a connection on a trivial bundle. For a general vector bundle $E$, 
    The existence of connections on $E$ can be proved by using partitions of unity, which we left as an exercise.
\end{remark}

\begin{proposition}
    Let $\nabla_1$ and $\nabla_2$ be two connections on $E$, then $\nabla_1 - \nabla_2$ is $C^\infty(M)$-linear in $s$, 
    thus defines a section $A\in \Omega^1(M,\mathrm{End}(E))$. 
    Conversely, for any $A\in \Omega^1(M,\mathrm{End}(E))$ and a connection $\nabla$ on $E$, $\nabla + A$ is also a connection on $E$. 
    \textcolor{red}{Consequently, the space of connections on $E$ is an affine space modeled on $\Omega^1(M,\mathrm{End}(E))$.}
\end{proposition}
\begin{proof}
    For $s\in C^\infty(M,E),f\in C^\infty(M)$, 
    \begin{equation*}
        (\nabla_1 - \nabla_2)(fs) = \md f\otimes s + f\nabla_1 s - \md f\otimes s - f\nabla_2 s = f(\nabla_1 - \nabla_2)(s),
    \end{equation*}
    thus $\nabla_1 - \nabla_2$ is $C^\infty(M)$-linear in $s$. 
    Conversely, for $A\in \Omega^1(M,\mathrm{End}(E))$ and a connection $\nabla$ on $E$,
    \begin{equation*}
        (\nabla + A)(fs) = \nabla(fs) + A(fs) = \md f\otimes s + f\nabla s + fAs = \md f\otimes s + f(\nabla + A)(s). \qedhere
    \end{equation*}
\end{proof}

\subsubsection*{Curvature of connections}
\begin{proposition}
    $\nabla^2: \Omega^*(M,E)\to\Omega^{*+2}(M,E)$ is an $\Omega^*(M)$-linear map, i.e., 
    for any $\alpha\in \Omega^*(M),\eta\in \Omega^*(M,E)$,
    \begin{equation*}
        \nabla^2(\alpha\wedge\eta) = \alpha\wedge\nabla^2 \eta.
    \end{equation*}
\end{proposition}
\begin{proof}
    For $\alpha\in \Omega^q(M),\eta\in \Omega^*(M,E)$, 
    \begin{align*}
        \nabla^2\big(\alpha\wedge\eta) 
        =& \nabla(\md\alpha\wedge\eta 
        + (-1)^{q}\alpha\wedge \nabla\eta\big) \\
        =& \md^2\alpha\wedge\eta + (-1)^{q+1}\md\alpha\wedge \nabla\eta 
        + (-1)^q \md\alpha\wedge \nabla\eta 
        + \alpha\wedge\nabla^2 \eta \\
        =& \alpha\wedge\nabla^2 \eta. \qedhere
    \end{align*}
\end{proof}
\begin{definition}
    The curvature of $\nabla$ is defined as $R:=\nabla^2\in\Omega^2(M,\mathrm{End}(E))$.
\end{definition}

\subsubsection*{Metrics and metric compatibility}
\begin{definition}
    Let $E$ be a (complex) vector bundle over $M$.
    A \textbf{(Hermitian) metric} on $E$ is a smooth section $g$ of $E^*\otimes E^*$ such that 
    for any $x\in M$, $g_x$ is a (Hermitian) inner product $\langle \cdot, \cdot \rangle_{E_x}$ on the fibre $E_x$. 
\end{definition}
\begin{definition}
    A connection $\nabla$ on $E$ is called \textbf{compatible with the metric $g$} 
    if for any $X\in C^\infty(M,TM)$ and $s_1,s_2\in C^\infty(M,E)$, 
    \begin{equation*}
        X\langle s_1, s_2 \rangle = \langle \nabla_X s_1, s_2 \rangle + \langle s_1, \nabla_X s_2 \rangle.
    \end{equation*}
\end{definition}
\begin{remark}
    For a (complex) vector bundle $E$ over $M$, we can always construct a (Hermitian) metric on $E$ by using partitions of unity. 
\end{remark}
\begin{remark}
    Given a (Hermitian) metric $g$ on $E$, we can always construct a connection $\nabla$ 
    compatible with $g$ by using partitions of unity. Note that the connection compatible with $g$ is not unique.
\end{remark}
\begin{proposition}
    Let $\nabla_1$ and $\nabla_2$ be two connections on $E$ compatible with metric $g$, 
    then $\nabla_1 - \nabla_2$ is skew-symmetric w.r.t. $g$, i.e., $\nabla_1 - \nabla_2 \in \Omega^1(M,\mathfrak{so}(E,g))$. 
    Conversely, for any $A\in \Omega^1(M,\mathfrak{so}(E,g))$ and a connection $\nabla$ compatible with $g$, 
    $\nabla + A$ is also compatible with $g$. 
    \textcolor{red}{Consequently, the space of connections on $E$ compatible with $g$ 
    is an affine space modeled on $\Omega^1(M,\mathfrak{so}(E,g))$.}
\end{proposition}
\begin{proof}
    For $s_1,s_2\in C^\infty(M,E), X\in C^\infty(M,TM)$, 
    \begin{gather*}
        \begin{cases}
            X\langle s_1, s_2 \rangle = \langle \nabla_{1,X} s_1, s_2 \rangle + \langle s_1, \nabla_{1,X} s_2 \rangle \\
            X\langle s_1, s_2 \rangle = \langle \nabla_{2,X} s_1, s_2 \rangle + \langle s_1, \nabla_{2,X} s_2 \rangle
        \end{cases} \\ 
        \implies \langle (\nabla_1 - \nabla_2)_X s_1, s_2 \rangle + \langle s_1, (\nabla_1 - \nabla_2)_X s_2 \rangle = 0.
    \end{gather*}
    Conversely, for $A\in \Omega^1(M,\mathfrak{so}(E,g))$ and a connection $\nabla$ compatible with $g$,
    \begin{align*}
        &\langle (\nabla + A)_X s_1, s_2 \rangle + \langle s_1, (\nabla + A)_X s_2 \rangle \\
        =& \langle \nabla_X s_1, s_2 \rangle + \langle s_1, \nabla_X s_2 \rangle + \langle A(X) s_1, s_2 \rangle + \langle s_1, A(X) s_2 \rangle \\
        =& X\langle s_1, s_2 \rangle. \qedhere
    \end{align*}
\end{proof}
% \begin{definition}
%     For $A\in C^\infty(M,\mathrm{End}(E))$, we define $\nabla A \in \Omega^{1}(M,\mathrm{End}(E))$ by 
%     \begin{equation*}
%         (\nabla A)(s) := \nabla(A(s)) - A(\nabla s).
%     \end{equation*}
%     this is a connection on $\mathrm{End}(E)$ induced by $\nabla$ on $E$. It can also be extended to 
%     $\nabla: \Omega^*(M,\mathrm{End}(E)) \to \Omega^{*+1}(M,\mathrm{End}(E))$ by the Leibniz rule:
%     \begin{equation*}
%         \nabla(\alpha\wedge A) = \md\alpha\wedge A + (-1)^{q}\alpha\wedge \nabla A,
%         \quad\forall \alpha\in \Omega^q(M),A\in \Omega^*(M,\mathrm{End}(E)).
%     \end{equation*}
% \end{definition}
\subsubsection{Connection 1-form}
Let $E$ be a rank $m$ vector bundle over $M$, and $\{e_i\}$ be a local frame of $E$ on an open set $U\subset M$. 
Define $\omega_j^k$ by
\begin{equation*}
    \nabla_X e_j = \omega_j^k(X) e_k,
\end{equation*}
this 1-form matrix $\omega = (\omega_j^k)$ is called \textbf{the connection 1-form} of $\nabla$ w.r.t. $\{e_i\}$. 
Then we can write the above equation in a more compact way:
\begin{equation*}
    \nabla (\ft{e}{1}{m}) = (\ft{e}{1}{m})\,\omega = 
    (\ft{e}{1}{m})
    \begin{pmatrix}
        \omega_1^1 & \cdots & \omega_1^m \\
        \vdots &  & \vdots \\
        \omega_m^1 & \cdots & \omega_m^m
    \end{pmatrix}.
\end{equation*}
Thus for any section $s = s^i e_i$, we have 
\begin{align*}
    \nabla s =& \nabla\left((\ft{e}{1}{m})
    \begin{pmatrix}
        s^1 \\ \vdots \\ s^m
    \end{pmatrix}
    \right) 
    = \Big(\nabla(\ft{e}{1}{m})\Big)
    \begin{pmatrix}
        s^1 \\ \vdots \\ s^m
    \end{pmatrix} + (\ft{e}{1}{m}) 
    \begin{pmatrix}
        \md s^1 \\ \vdots \\ \md s^m
    \end{pmatrix} \\
    =& (\ft{e}{1}{m}) 
    \left( \md 
    \begin{pmatrix}
        s^1 \\ \vdots \\ s^m 
    \end{pmatrix} + 
    \omega 
    \begin{pmatrix}
        s^1 \\ \vdots \\ s^m
    \end{pmatrix}
    \right).
\end{align*}
So under a local frame, $\nabla$ can be expressed as
\begin{equation*}
    \nabla = \md + \omega.
\end{equation*}
\begin{proposition}
    Suppose $\{\tilde{e}_i\}$ are another local frame of $E$ on $U$, 
    and the transition matrix from $\{e_i\}$ to $\{\tilde{e}_i\}$ is $a=(a_i^j)$, i.e.
    \begin{equation*}
        (\ft{\tilde{e}}{1}{m}) = (\ft{e}{1}{m})\, a,
    \end{equation*}
    Let $\tilde{\omega}$ be the connection 1-form w.r.t. $\{\tilde{e}_i\}$,
    then the connection 1-forms transform as
    \begin{equation*}
        \tilde{\omega} = a^{-1} \md a + a^{-1} \omega a.
    \end{equation*}
\end{proposition}

A local frame of $E$ is the same as a local trivialization. 
Let $\{(U_\alpha,\psi_\alpha)\}$ be a local trivialization of $E$, $g_{\alpha\beta}$ be the transition functions. 
A section $s\in C^\infty(M,E)$ possesses a local representative $s_\alpha$, which is an $m$-tuple of smooth functions on $U_\alpha$. 
Also, $\nabla s$ possesses a local representative $(\nabla s)_\alpha$, which is an $m$-tuple of 1-forms on $U_\alpha$. 
Just as in the preceding paragraph, we can write 
\begin{equation}\label{eq:local-expression-of-connection}
    (\nabla s)_\alpha = \md s_\alpha + \omega_\alpha s_\alpha,
\end{equation}
where $\omega_\alpha$ is an $m\times m$ matrix of 1-forms on $U_\alpha$. 

Since the connection is globally defined, the local representatives of $\nabla s$ must satisfy
\begin{equation*}
    \md s_\beta + \omega_\beta s_\beta = g_{\alpha\beta}^{-1}(\md s_\alpha + \omega_\alpha s_\alpha) 
    \qquad \text{on }\; U_\alpha\cap U_\beta.
\end{equation*}
Apply $s_\alpha = g_{\alpha\beta} s_\beta$ to the above equation, 
\begin{align*}
    \md s_\beta + \omega_\beta s_\beta 
    =& g_{\alpha\beta}^{-1}\big(\md (g_{\alpha\beta} s_\beta) + \omega_\alpha g_{\alpha\beta} s_\beta\big) \\
    =& g_{\alpha\beta}^{-1}\big((\md g_{\alpha\beta}) s_\beta + g_{\alpha\beta} \md s_\beta + \omega_\alpha g_{\alpha\beta} s_\beta\big) \\
    =& \md s_\beta + \Big(g_{\alpha\beta}^{-1}\md g_{\alpha\beta} + g_{\alpha\beta}^{-1}\omega_\alpha g_{\alpha\beta}\Big) s_\beta,    
\end{align*}
Thus 
\begin{equation}\label{eq:translation-formula-of-connection-1-forms}
    \omega_\beta = g_{\alpha\beta}^{-1}\md g_{\alpha\beta} + g_{\alpha\beta}^{-1}\omega_\alpha g_{\alpha\beta}.
\end{equation}
The above discussion can be summarized as the following theorem:
\begin{theorem}
    Let $\{(U_\alpha,\psi_\alpha)\}$ be a local trivialization of $E$, with transition functions $g_{\alpha\beta}$, 
    and $\omega_\alpha$ be an $m\times m$ matrix of 1-forms on $U_\alpha$ 
    satisfying the transformation formula {\rm(}\ref{eq:translation-formula-of-connection-1-forms}{\rm)}. 
    This imformation uniquely determines a connection $\nabla$ on $E$. 
\end{theorem}

Let $(E,g)$ be a real vector bundle equipped with a metric, we can choose the local frame $\{e_i\}$ on $U$ to be orthonormal, 
then for any $X\in C^{\infty}(U, TU)$, 
\begin{equation*}
    0=X\langle e_i, e_j\rangle=\langle\omega_i^k(X)e_k, e_j\rangle + \langle e_i, \omega_j^k(X)e_k\rangle = \omega_j^i(X) + \omega_i^j(X),
\end{equation*}
so the 1-form matrix $\omega$ is skew-symmetric, i.e., $\omega \in \Omega^1(U,\mathfrak{so}(m))$. 

Let $\{(U_\alpha,\psi_\alpha)\}$ be a local trivialization of $(E,g)$. We can assume 
\begin{equation*}
    \psi_\alpha: \pi^{-1}(U_\alpha) \to U_\alpha \times \mathbb{R}^m
\end{equation*}
is isometric on each fibre, where the metric on $U_\alpha\times\mathbb{R}^m$ is Euclidean metric on each fibre. 
Then the transition functions $g_{\alpha\beta}$ take values in $O(m)$, 
and each $\omega_\alpha$ is in $\Omega^1(U_\alpha,\mathfrak{so}(m))$. 

If $E$ is a complex vector bundle and $g$ is a Hermitian metric. Then all the discussion above still holds, 
except that $\omega_\alpha$ is in $\Omega^1(U_\alpha,\mathfrak{u}(m))$ and $g_{\alpha\beta}$ takes values in $U(m)$. 

\begin{example}
    Let $\pi:L\to M$ be a complex line bundle, $g$ be a Hermition metric on $L$ and $\nabla$ be a connection on $L$ compatible with $g$. 
    Choose a local trivialization $\{(U_\alpha,\psi_\alpha)\}$ of $L$ such that $\psi_\alpha: \pi^{-1}(U_\alpha) \to U_\alpha \times \mathbb{C}$ is an isometry. 
    Then $g_{\alpha\beta}$ takes values in $U(1)$, so we can write $g_{\alpha\beta} = e^{i\theta_{\alpha\beta}}$ 
    for some real-valued function $\theta_{\alpha\beta}$ on $U_\alpha\cap U_\beta$. Meanwhile, $\omega_\alpha$ takes values in $\mathfrak{u}(1)$, 
    so we can write $\omega_\alpha = i A_\alpha$ for some real-valued 1-form $A_\alpha$ on $U_\alpha$.
    For any section $s$ of $L$, on $U_\alpha$, we have
    \begin{equation*}
        (\nabla s)_\alpha = \md s_\alpha + \omega_\alpha s_\alpha, \qquad \omega_\alpha = i A_\alpha,
    \end{equation*}
    On $U_\alpha\cap U_\beta$, 
    \begin{equation*}
        \omega_\beta = g_{\alpha\beta}^{-1}\md g_{\alpha\beta} + g_{\alpha\beta}^{-1}\omega_\alpha g_{\alpha\beta} 
        = e^{-i\theta_{\alpha\beta}} \md e^{i\theta_{\alpha\beta}} + e^{-i\theta_{\alpha\beta}} \omega_\alpha e^{i\theta_{\alpha\beta}} 
        = \omega_\alpha + i\md \theta_{\alpha\beta} \Big(= \omega_\alpha - i\md \theta_{\beta\alpha}\Big)
    \end{equation*}
    Recall $U(1)$-gauge theory in physics, 
    the following table shows the correspondence between concepts in physics and differential geometry:
    \begin{center}
    \begin{tabular}{|c|c|}
        \hline
        \textbf{Physics} & \textbf{Differential Geometry} \\
        \hline
        wave function $\psi$ & section $s$ of complex line bundle $L$ \\
        \hline
        $\psi\mapsto\psi'=e^{i\chi}\psi$ & $s_\beta=g_{\beta\alpha}s_\alpha$ \\
        \hline
        covariant derivative & connection $\nabla$ on $L$ \\
        \hline
        $D'\psi' = e^{i\chi} D\psi$ & $(\nabla s)_\beta=g_{\beta\alpha}(\nabla s)_\alpha$ \footnotemark \\
        \hline
        $D=\md+iA$ & $(\nabla s)_\alpha=\md s_\alpha+\omega_\alpha s_\alpha$ \\
        \hline
        $A'=A-\md\chi$ & $\omega_\beta=\omega_\alpha-i\md\theta_{\beta\alpha}$ \\
        \hline
    \end{tabular}
\end{center}
\footnotetext{That is to say, $\nabla s$ is globally defined.}
\end{example}

\subsubsection{Horizontal Distribution}
In this subsection, we give another interpretation of a connection on a vector bundle, namely a horizontal distribution, 
and prove that it is equivalent to the definition given in terms of covariant derivatives.

Let $\pi:E\to M$ be a rank $m$ vector bundle, and $p\in E_x$ be a point in the fibre over $x\in M$. 
The propjection $\pi$ induces a linear map $\pi_*{}_p: T_p E \to T_x M$. 
We define $\mathcal{V}_p:=\ker \pi_*{}_p$, which is called the \textbf{vertical tangent space} at $p$. 
Thus we have a short exact sequence 
\[\begin{tikzcd}
    0 & {\mathcal{V}_p} & T_pE & {T_{\pi(p)}M} & 0
    \arrow[from=1-1, to=1-2]
    \arrow[ hook, from=1-2, to=1-3]
    \arrow["\pi_*", from=1-3, to=1-4]
    \arrow[from=1-4, to=1-5]
\end{tikzcd}\]
and
\begin{equation*}
    \dim \mathcal{V}_p = \dim T_p E - \dim T_{\pi(p)} M = m.
\end{equation*}
The kernel of $\pi_*$ form a subbundle of $TE$. 
We call it the \textbf{vertical subbundle} of $TE$ and denote it by $\mathcal{V}$.
\begin{proposition}
    For each $p\in E$, there is an isomorphism $E_{\pi(p)} \cong \mathcal{V}_p$, which induces an 
    isomorphism $\pi^* E \cong \mathcal{V}$. 
\end{proposition}
\begin{proof}
    A point $(p,p')\in \pi^* E$ must satisfy $\pi(p) = \pi(p')$, so we can define a curve $\gamma(t) = p + t p'$ in $E$.
    Then $\gamma'(0) \in T_p E$ gives a tangent vector at $p$, this gives an injective morphism 
    \begin{equation*}
        \pi^* E \hookrightarrow TE, \quad (p,p') \mapsto (p,\gamma'(0)).
    \end{equation*}
    So we can view $\pi^* E$ as a subbundle of $TE$. 
    Moreover, since $\gamma(t)\subset E_{\pi(p)}$, we have $\pi(\gamma(t)) = \pi(p)$, thus $\pi_*{}_p(\gamma'(0)) = 0$. 
    Therefore, $\pi^* E\subset \ker\pi_*$ 

    Since $\pi_*{}_p$ is a surjective linear map, we have 
    \begin{equation*}
        \dim \ker \pi_*{}_p = \dim T_p E - \dim T_{\pi(p)} M = m = (\pi^* E)_p.
    \end{equation*}
    Thus $\pi^* E = \ker \pi_*$.
\end{proof}
For any $p\in E$, we can define a map 
\begin{equation*}
    \iota_p: E_{\pi(p)} \to \mathcal{V}_p\subset T_p E, \quad p' \mapsto \gamma'(0),
\end{equation*}
where $\gamma(t) = p + t p'$ is a curve in $E$. 
Then $\iota_p$ is an isomorphism between $E_{\pi(p)}$ and $\mathcal{V}_p$.

Let $\pi^*TM$ be the pull-back of $TM$ by $\pi$, then $\pi_*: TE \to TM$ can be viewed as 
a morphism $\tilde{\pi}_*: TE \to \pi^* TM$. 
The above proposition implies the following short exact sequence of vector bundles over $E$:
% \begin{equation*}
%     0 \to \pi^*E \overset{i}{\hookrightarrow} TE \xrightarrow{\tilde{\pi}_*} \pi^* TM \to 0.
% \end{equation*}
\[\begin{tikzcd}
    0 & \mathcal{V} & TE & {\pi^*TM} & 0
    \arrow[from=1-1, to=1-2]
    \arrow[hook, from=1-2, to=1-3]
    \arrow["\tilde{\pi}_*", from=1-3, to=1-4]
    \arrow[from=1-4, to=1-5]
\end{tikzcd}\]
In smooth category, all short exact sequences of vector bundles split, 
so we can find a subbundle $\mathcal{H}\hookrightarrow TE$ such that 
\begin{equation*}
    TE = \mathcal{H} \oplus \mathcal{V}, \qquad \tilde{\pi}_*: \mathcal{H} \to \pi^* TM \text{ is an isomorphism}.
\end{equation*}

On the vector bundle $E$, we have two natural operations: scalar multiplication and addition. 
\begin{itemize}
    \item The scalar multiplication: 
    \begin{equation*}
        m: \mathbb{R} \times E \to E, \quad m(\lambda,p) = \lambda p.
    \end{equation*}
    When fixing $\lambda$, we get a map $m_\lambda: E \to E,\, m_\lambda(p) = \lambda p$, 
    which is the scalar multiplication by $\lambda$.
    % For $\nu\in T_\lambda \mathbb{R}\cong \mathbb{R}$ and $v\in T_p E$, let $\gamma(t)$ be a curve in $E$ such that 
    % $\gamma(0) = p$ and $\gamma'(0) = v$, then 
    % \begin{align*}
    %     m_*{}_{(\lambda,p)}(\nu,v) =& \frac{\md }{\md t}\Big|_{t=0} (\lambda + t\nu)\gamma(t) \\
    %     =& \frac{\md }{\md t}\Big|_{t=0}(\lambda p + t\nu p) + \frac{\md }{\md t}\Big|_{t=0} \lambda \gamma(t) \\
    %     =& \nu\iota_p(p) + m_\lambda{}_*{}_p v,
    % \end{align*} 
    \item The addition:
    \begin{equation*}
        a: E\times_M E \to E, \quad a(p,q) = p + q,
    \end{equation*}
    where $E\times_M E = \{(p,q)\in E\times E: \pi(p) = \pi(q)\}$. 
    % For $u\in T_p E$ and $v\in T_q E$ such that $\pi(p) = \pi(q)$, 
    % let $\gamma(t)$ and $\sigma(t)$ be curves in $E$ such that $\gamma(0) = p$, $\gamma'(0) = u$, $\sigma(0) = q$ and $\sigma'(0) = v$, 
    % then 
    % \begin{align*}
    %     a_*{}_{(p,q)}(u,v) =& \frac{\md }{\md t}\Big|_{t=0} a(\gamma(t), \sigma(t)) \\
    %     =& \frac{\md }{\md t}\Big|_{t=0} a(\gamma(t), q) + \frac{\md }{\md t}\Big|_{t=0} a(p, \sigma(t)) \\
    %     =& a_q{}_*{}_p u + a_p{}_*{}_q v,
    % \end{align*}
    % where $a_q: E_{\pi(p)} \to E_{\pi(p)},\, a_q(p) = p + q$ is the addition by $q$, 
    % and $a_p: E_{\pi(q)} \to E_{\pi(q)},\, a_p(q) = p + q$ is the addition by $p$.
\end{itemize}

\begin{definition}
    A \textbf{horizontal distribution} of $E$ is a smooth subbundle $\mathcal{H}\subset TE$ 
    satisfying the following properties 
    \begin{enumerate}
        \item There exists a splitting map $k: \pi^* TM \to TE$ such that $\tilde{\pi}_* \circ k = \mathrm{id}$.
        \[\begin{tikzcd}
            0 & \mathcal{V} & TE & {\pi^*TM} & 0
            \arrow[from=1-1, to=1-2]
            \arrow[hook, from=1-2, to=1-3]
            \arrow["\tilde{\pi}_*", shift left, from=1-3, to=1-4]
            \arrow["k", shift left, from=1-4, to=1-3]
            \arrow[from=1-4, to=1-5]
        \end{tikzcd}\]
        That is to say, 
        \begin{flalign*}
            &\text{(H1)} && TE=\mathcal{H}\oplus\mathcal{V}. &
        \end{flalign*}
        \item $\mathcal{H}$ is invariant under the scalar multiplication on $E$, i.e. for any $\lambda\in \mathbb{R},\,p\in E$,
        \begin{flalign*}
            &\text{(H2)} && m_{\lambda}{}_*{}_p \mathcal{H}_p = \mathcal{H}_{\lambda p}. &
        \end{flalign*}
        \item $\mathcal{H}$ is invariant under the addition on $E$, i.e. for any $(p,q)\in E\times_M E$,
        \begin{flalign*}
            &\text{(H3)} && a_*{}_{(p,q)}(\mathcal{H}_p,\mathcal{H}_q) = \mathcal{H}_{p+q}. &
        \end{flalign*} 
    \end{enumerate}
    The last two conditions are saying that $\mathcal{H}$ is compatible with the linear structure of $E$.
\end{definition}

There is no canonical way to choose a horizontal distribution, so a horizontal distribution is an extra structure on $E$.

In what follows, we reveal that covariant derivatives and horizontal distributions can actually determine each other.

\subsubsection*{Covariant Derivative determines Horizontal Distribution}
Let $\nabla$ be a connection on $E$. Given any smmoth curve $\gamma: (-\varepsilon,\varepsilon) \to M$ with $\gamma(0) = x$, 
we say $\tilde{\gamma}: (-\varepsilon,\varepsilon) \to E$ is a 
\textbf{lift} of $\gamma$ if $\pi(\tilde{\gamma}(t)) = \gamma(t)$ for all $t$. 
\begin{definition}
    A lift $\tilde{\gamma}$ of $\gamma$ is called a \textbf{horizontal lift} if 
    \begin{equation*}
        \nabla_{\gamma'(t)} \tilde{\gamma}(t) = 0, \quad \forall t\in (-\varepsilon,\varepsilon).
    \end{equation*}
\end{definition}
\begin{example}
    In the case of a Riemannian manifold $(M,g)$, the horizontal lift is exactly the parallel transport.
\end{example}
\begin{proposition}
    For any (picewise) smooth curve $\gamma: (-\varepsilon,\varepsilon) \to M$ and any point $p\in E_{\gamma(0)}$, 
    there exists a unique horizontal lift $\tilde{\gamma}$ of $\gamma$ starting at $p$.
\end{proposition}
\begin{proof}[Sketch of proof]
    The proof is exactly the same as the existence and uniqueness of geodesics. 
    The outline of the proof is to write the equations as a system of ODEs in a local frame; 
    the conclusion then follows from the existence and uniqueness theory of linear ODEs.
\end{proof}
For any vector $u\in T_x M$ and any point $p\in E_x$, 
we can find a curve $\gamma$ such that $\gamma(0)=x$ and $\gamma'(0)=u$, 
the horizontal lift $\tilde{\gamma}$ gives a tangent vector $\tilde{\gamma}'(0)\in T_p E$. 
\begin{fact}
    $\tilde{\gamma}'(0)\in T_pE$ is independent of the choice of $\gamma$.
\end{fact}
\begin{proof}[Sketch of proof]
    By solving the ODEs, we know that $\tilde{\gamma}'(0)$ depends only on the derivative of $\gamma$ at $0$, 
    which is $\gamma'(0)=u$.
\end{proof}
So there is a well-defined map 
\begin{equation*}
    T_{\pi(p)} M \to T_p E, \quad (u^h)_p := \tilde{\gamma}'(0).
\end{equation*}
We call $(u^h)_p$ the \textbf{horizontal lift} of $u$ at $p$.
By definition, 
\begin{equation*}
    \pi(\tilde{\gamma}(t)) = \gamma(t) \implies 
    \pi_*{}_p(u^h)_p = \pi_*{}_p(\tilde{\gamma}'(0)) = \gamma'(0) = u.
\end{equation*}
So we obtain a linear injection from $T_{\pi(p)} M$ to $T_p E$, 
whose image forms a distribution on $E$ denoted by $\mathcal{H}$, i.e., 
\begin{equation*}
    \textcolor{red}{\mathcal{H} := \big\{(p,(u^h)_p)\mid p\in E, u\in T_{\pi(p)} M\big\}.}
\end{equation*}
\begin{fact}
    $\mathcal{H}$ is a smooth subbundle of $TE$. 
\end{fact}
\begin{proof}[Sketch of proof]
    This fact relies on the smooth dependence of solutions of linear ODEs on initial conditions and parameters.
\end{proof}
\begin{fact}
    $\mathcal{H}$ is invariant under the scalar multiplication and addition on $E$.
\end{fact}
\begin{proof}
    Let $p,q\in E$ such that $\pi(p) = \pi(q)$, $\lambda\in \mathbb{R}$, and $\gamma$ be a curve such that $\gamma(0) = \pi(p)$. 
    Let $\tilde{\gamma}_p,\tilde{\gamma}_q$ be the horizontal lift of $\gamma$ starting at $p,q$ respectively, 
    then by the uniqueness of horizontal lift, $\lambda \tilde{\gamma}_p$ is the horizontal lift of $\gamma$ starting at $\lambda p$,
    and $\tilde{\gamma}_p + \tilde{\gamma}_q$ is the horizontal lift of $\gamma$ starting at $p+q$. Thus 
    \begin{gather*}
        m_{\lambda}{}_*{}_p \tilde{\gamma}'(0) = (\lambda \tilde{\gamma}_p)'(0) \in \mathcal{H}_{\lambda p}, \\
        a_*{}_{(p,q)}(\tilde{\gamma}_p'(0), \tilde{\gamma}_q'(0)) = (\tilde{\gamma}_p + \tilde{\gamma}_q)'(0) \in \mathcal{H}_{p+q}. \qedhere
    \end{gather*}
\end{proof}
The above discussion is summarized in the following theorem.
\begin{theorem}
    $\mathcal{H}$ is a horizontal distribution of $E$. 
\end{theorem}

% \begin{example}
%     For a Riemannian manifold $(M,g)$ and the Levi-Civita connection $\nabla$, 
%     the tangent bundle $T(TM)$ has a decomposition $T(TM) = \mathcal{V}\oplus\mathcal{H}$. 
%     We can define a metric $G$ on $TM$ by
%     \begin{gather*}
%         G(X^h,Y^h):= g(\pi_*(X^h), \pi_*(Y^h)), \quad G(X^v,Y^v) := g(\iota^{-1}(X^v), \iota^{-1}(Y^v)), \\
%         \quad G(X^h,Y^v) := 0,
%     \end{gather*} 
%     where $X^h,Y^h\in \mathcal{H}$ and $X^v,Y^v\in \mathcal{V}$. 
%     $G$ is called \textbf{the Sasaki metric} on $TM$.
% \end{example}

\subsubsection*{Horizontal Distribution determines Covariant Derivative}
Conversaly, given a horizontal distribution $\mathcal{H}$, we can derive a connection $\nabla$ on $E$ as follows. 
\begin{definition}
    For any $p\in E$ and $u\in TE$, we can write $u = h(u) + v(u)$ where 
    $h(u)\in \mathcal{H}$ and $v(u)\in \mathcal{V}$. We call $h(u)$ and $v(u)$ 
    the horizontal and vertical component of $u$ respectively.
\end{definition}
For any $X\in C^\infty(M,TM)$ and $s\in C^\infty(M,E)$, the tangent map of $s$ gives a map $s_*: TM \to TE$, 
then at any point $x\in M$ we can define $(\nabla_X s)(x)$ by
\begin{equation*}
    \textcolor{red}{(\nabla_X s)(x) := \iota_{s(x)}^{-1}\big(v(s_*{}_x X_x)\big).}
\end{equation*}
Note that $\iota_{s(x)}: E_x \to \mathcal{V}_{s(x)}$ is the isomorphism defined above. 
The idea is taking the vertical component of $s_*{}_x(X_x)$.
\begin{theorem}\label{thm:connection-from-horizontal-distribution}
    $\nabla$ defined above is a connection on $E$.
\end{theorem}
To prove this theorem, we need the following lemma.
\begin{lemma}\label{lem:tangent-map-of-linear-ops-acts-on-vertical-bundle}
    For $\lambda\in\mathbb{R},\,(p,q)\in E\times_M E$ and $u,w\in E_{\pi(p)}$, we have  
    \begin{gather*}
        m_{\lambda}{}_*{}_p \iota_p(u) = \iota_{\lambda p}(\lambda u), \\
        a_*{}_{(p,q)}(\iota_p(u),\iota_q(w)) = \iota_{p+q}(u+w).
    \end{gather*}
\end{lemma}
\begin{proof}
    By the definition of $\iota_p$, 
    \begin{gather*}
        m_{\lambda}{}_*{}_p \iota_p(u) = \left.\frac{\md }{\md t}\right|_{t=0} m_{\lambda}(p + t u) 
        = \left.\frac{\md }{\md t}\right|_{t=0} (\lambda p + t \lambda u) = \iota_{\lambda p}(\lambda u), \\
        a_*{}_{(p,q)}(\iota_p(u),\iota_q(w)) = \left.\frac{\md }{\md t}\right|_{t=0} (p + t u) + (q + t w) 
        = \iota_{p+q}(u+w). \qedhere
    \end{gather*}
\end{proof}

\begin{lemma}\label{lem:H-compatible-with-linear-str}
    For any $\lambda\in\mathbb{R},\,(p,q)\in E\times_M E$ and $u\in TE_p,\,w\in TE_q$, 
    \begin{gather*}
        h(m_{\lambda}{}_*{}_p u) = m_{\lambda}{}_*{}_p(h(u)), \quad v(m_{\lambda}{}_*{}_p u) = m_{\lambda}{}_*{}_p(v(u)), \\
        h(a_*{}_{(p,q)}(u,w)) = a_*{}_{(p,q)}(h(u),h(w)), \quad v(a_*{}_{(p,q)}(u,w)) = a_*{}_{(p,q)}(v(u),v(w)).
    \end{gather*}
\end{lemma}
\begin{proof}
    Decompose $u$ and $m_{\lambda}{}_*{}_p u$ into their horizontal and vertical components, we get 
    \begin{equation*}
        m_{\lambda}{}_*{}_p(h(u)) + m_{\lambda}{}_*{}_p(v(u)) = m_{\lambda}{}_*{}_p(u) 
        = h(m_{\lambda}{}_*{}_p u) + v(m_{\lambda}{}_*{}_p u).
    \end{equation*}
    Since $\mathcal{H}$ is compatible with the linear structure of the vector bundle, 
    we have $m_{\lambda}{}_*{}_p(h(u))\in \mathcal{H}_{\lambda p}$. Besides, 
    Lemma \ref{lem:tangent-map-of-linear-ops-acts-on-vertical-bundle} implies 
    $m_{\lambda}{}_*{}_p(v(u))\in \mathcal{V}_{\lambda p}$, thus 
    \begin{equation*}
        m_{\lambda}{}_*{}_p(h(u)) = h(m_{\lambda}{}_*{}_p u), \qquad m_{\lambda}{}_*{}_p(v(u)) = v(m_{\lambda}{}_*{}_p u).
    \end{equation*}
    The proof for the addition is exactly the same.
\end{proof}

\begin{lemma}\label{lem:tangent-map-of-multiplication}
    For $\nu\in T_\lambda \mathbb{R}\cong \mathbb{R}$ and $u\in T_p E$, we have 
    \begin{equation*}
        m_*{}_{(\lambda,p)}(\nu,u) = \nu\cdot\iota_p(p) + m_\lambda{}_*{}_p u.
    \end{equation*}
\end{lemma}
\begin{proof}
    let $\gamma(t)$ be a curve in $E$ such that $\gamma(0) = p$ and $\gamma'(0) = u$, then 
    \begin{align*}
        m_*{}_{(\lambda,p)}(\nu,u) =& \left.\frac{\md }{\md t}\right|_{t=0} (\lambda + t\nu)\gamma(t) \\
        =& \left.\frac{\md }{\md t}\right|_{t=0}(\lambda p + t\nu p) 
        + \left.\frac{\md }{\md t}\right|_{t=0} \lambda \gamma(t) \\
        =& \nu\cdot\iota_p(p) + m_\lambda{}_*{}_p u. \qedhere
    \end{align*} 
\end{proof}
\begin{proof}[Proof of Theorem \ref{thm:connection-from-horizontal-distribution}]
    We need to verify that $\nabla$ satisfies the properties of a connection. 
    \begin{itemize}
        \item $\nabla_{f_1X_1 + f_2X_2} s = f_1 \nabla_{X_1} s + f_2 \nabla_{X_2} s$:
        \begin{align*}
            &\nabla_{f_1X_1 + f_2X_2} s(x) = \iota_{s(x)}^{-1}\Big(v\big(s_*{}_x(f_1X_1 + f_2X_2)\big)\Big) \\
            =& \iota_{s(x)}^{-1}\Big(v\big(s_*{}_x(f_1(x)X_1{}_x)\big) 
            + v\big(s_*{}_x(f_2(x)X_2{}_x)\big)\Big) \\
            =& \iota_{s(x)}^{-1}\Big(f_1(x) v\big(s_*{}_x X_1{}_x\big) 
            + f_2(x) v\big(s_*{}_x X_2{}_x\big)\Big) \\
            =& f_1 \nabla_{X_1} s(x) + f_2 \nabla_{X_2} s(x).
        \end{align*}
        \item $\nabla_X (\lambda_1 s_1 + \lambda_2 s_2) = \lambda_1 \nabla_X s_1 + \lambda_2 \nabla_X s_2$:
        
        For $\mathbb{R}$-linearity, 
        \begin{align*}
            &\nabla_X (\lambda s)(x) = \iota_{(\lambda s)(x)}^{-1}\Big(v\big((\lambda s)_*{}_x(X_x)\big)\Big) \\
            =& \iota_{(\lambda s)(x)}^{-1}\Big(v\big(m_{\lambda}{}_*{}_{s(x)}s_*{}_x(X_x)\big)\Big) \\
            =& \iota_{(\lambda s)(x)}^{-1}\Big(m_{\lambda}{}_*{}_{s(x)}v\big(s_*{}_x(X_x)\big)\Big) 
            \quad\text{(By Lemma \ref{lem:H-compatible-with-linear-str})}\\
            =& \iota_{(\lambda s)(x)}^{-1}\Big(\iota_{(\lambda s)(x)}(\lambda \nabla_X s(x))\Big) 
            \quad\text{(By Lemma \ref{lem:tangent-map-of-linear-ops-acts-on-vertical-bundle})}\\
            =& \lambda \nabla_X s(x),
        \end{align*}
        For additivity,
        \begin{align*}
            &\nabla_X (s_1 + s_2)(x) = \iota_{(s_1 + s_2)(x)}^{-1}\Big(v\big((s_1 + s_2)_*{}_x X_x\big)\Big) \\
            =& \iota_{(s_1 + s_2)(x)}^{-1}\Big(v\big(a_*{}_{(s_1(x),s_2(x))}(s_{1*}{}_x X_x, s_{2*}{}_x X_x)\big)\Big) \\
            =& \iota_{(s_1 + s_2)(x)}^{-1}\Big(a_*{}_{(s_1(x),s_2(x))}\big(v(s_{1*}{}_x X_x), v(s_{2*}{}_x X_x)\big)\Big) 
            \quad\text{(By Lemma \ref{lem:H-compatible-with-linear-str})}\\
            =& \iota_{(s_1 + s_2)(x)}^{-1}\Big(\iota_{(s_1 + s_2)(x)}(\nabla_X s_1(x) + \nabla_X s_2(x))\Big) 
            \quad\text{(By Lemma \ref{lem:tangent-map-of-linear-ops-acts-on-vertical-bundle})}\\
            =& \nabla_X s_1(x) + \nabla_X s_2(x).
        \end{align*}
        \item $\nabla_X (fs) = \md f(X)\cdot s + f\nabla_X s$:
        \begin{align*}
            &\nabla_X (fs)(x) = \iota_{(fs)(x)}^{-1}\Big(v\big((fs)_*{}_x X_x\big)\Big) \\
            =& \iota_{(fs)(x)}^{-1}\Big(v\big(m_{f(x)}{}_*{}_{s(x)}(f_*{}_x X_x, s_*{}_x X_x)\big)\Big) \\ 
            =& \iota_{(fs)(x)}^{-1}\Big(\md f(X)_x\cdot \iota_{(fs)(x)}(s(x)) + v\big(m_{f(x)}{}_*{}_{s(x)}s_*{}_x X_x\big)\Big) 
            \quad\text{(By Lemma \ref{lem:tangent-map-of-multiplication})}\\
            =& \iota_{(fs)(x)}^{-1}\Big(\iota_{(fs)(x)}\big(\md f(X)_x\cdot s(x) + f(x)\nabla_X s(x)\big)\Big) 
            \quad\text{(By Lemma \ref{lem:tangent-map-of-linear-ops-acts-on-vertical-bundle})}\\
            =& \md f(X)_x\cdot s(x) + f(x)\nabla_X s(x). \qedhere
        \end{align*}
    \end{itemize}
\end{proof}
% \subsubsection{Curvature of Connection}

\subsubsection{Parallel Transport}
Let $\nabla$ be a connection on $E$. For any (piecewise) smooth curve $\gamma: [a,b] \to M$ and any point $v\in E_{\gamma(a)}$, 
There exists a unique horizontal lift $\tilde{\gamma}_{v}$ of $\gamma$ starting at $v$. 
Varying $v\in E_{\gamma(a)}$, we get different point $w=\tilde{\gamma}_{v}(b)\in E_{\gamma(b)}$. 
This gives a map
\begin{equation*}
    P^\gamma: E_{\gamma(a)} \to E_{\gamma(b)}, \quad P^\gamma(v) := \tilde{\gamma_{v}}(b).
\end{equation*}
\begin{proposition}
    $P^\gamma$ is a linear isomorphism.
\end{proposition}
\begin{proof}
    For any $v,w\in E_{\gamma(a)}$ and $\lambda\in \mathbb{R}$, 
    the uniqueness of horizontal lift implies that 
    \begin{equation*}
        \lambda \tilde{\gamma}_v = \tilde{\gamma}_{\lambda v}, \qquad 
        \tilde{\gamma}_v + \tilde{\gamma}_w = \tilde{\gamma}_{v+w}.
    \end{equation*}
    Thus 
    \begin{gather*}
        P^\gamma(\lambda v) = \tilde{\gamma}_{\lambda v}(b) = (\lambda \tilde{\gamma}_v)(b) = \lambda P^\gamma(v), \\
        P^\gamma(v+w) = \tilde{\gamma}_{v+w}(b) = (\tilde{\gamma}_v + \tilde{\gamma}_w)(b) = P^\gamma(v) + P^\gamma(w).
    \end{gather*}
    So $P^\gamma$ is a linear map. 
    Let $\bar{\gamma}(t) := \gamma(a+b-t)$ be the reversed curve of $\gamma$, 
    then one can check that $P^{\bar{\gamma}}: E_{\gamma(b)} \to E_{\gamma(a)}$ is the inverse map of $P^\gamma$. 
\end{proof}
\begin{definition}
    $P^\gamma$ is called \textbf{the parallel transport} along $\gamma$.
\end{definition}
\begin{remark}
    If $E$ possesses a metric $g$ and $\nabla$ is compatible with $g$, 
    then $P^\gamma: E_{\gamma(a)} \to E_{\gamma(b)}$ is an isometry. 
    The converse is also true: 
    if $P^\gamma$ is an isometry for any curve $\gamma$, 
    then $\nabla$ is compatible with $g$.
\end{remark}
\begin{remark}
    $P^\gamma$ highly depends on the choice of $\gamma$, not only on the endpoints $\gamma(a)$ and $\gamma(b)$.
    In general, if $\gamma_1$ and $\gamma_2$ are two curves with the same endpoints, 
    $P^{\gamma_1}$ and $P^{\gamma_2}$ can be very different. 
    In other words, if $\gamma$ is a loop, 
    $P^\gamma: E_{\gamma(a)} \to E_{\gamma(a)}$ may not be the identity map.
    % We can define the \textbf{holonomy group} of $\nabla$ at $x\in M$ as the subgroup of $\mathrm{GL}(E_x)$ consisting of all $P^\gamma$ where $\gamma$ is a loop based at $x$.
\end{remark}

Given a connection $\nabla$ on $E$, for any two points $x,y\in M$, 
we can link $x$ and $y$ by a curve $\gamma$, 
then $\nabla$ gives a linear isomorphism $P^\gamma: E_x \to E_y$. 
In this sense, $\nabla$ \textcolor{red}{connects} any two fibers of $E$.
This is the reason why $\nabla$ is called a \textbf{connection}. 

The key point is that if any curve $\gamma$ linking $x$ and $y$ gives an isomorphism $\Phi^\gamma: E_x \to E_y$ 
and $\Phi^\gamma$ satisfies some natural properties, then we can define a connection $\nabla$ on $E$ such that $P^\gamma = \Phi^\gamma$.
So intuitively, \textcolor{red}{a connection can be thought of as a family of isomorphisms between fibers at endpoints of curves}.

\begin{proposition}
    Let $\nabla$ be a connection on $E$, $X_{x}\in T_xM$ be a tangent vector at $x$ 
    and $s:M\to E$ be a section of $E$. Then there exists a curve $\gamma: (-\varepsilon,\varepsilon) \to M$ such that $\gamma(0) = x$ and $\gamma'(0) = X_x$. 
    If we denote $P_t: E_x \to E_{\gamma(t)}$ as the parallel transport along $\gamma|_{[0,t]}$, then 
    \begin{equation}\label{eq:con-def-by-para}
        \nabla_{X_x} s = \left.\frac{\md }{\md t}\right|_{t=0} P_t{}^{-1}(s(\gamma(t))).
    \end{equation}
\end{proposition}
\begin{proof}
    Let $\ft{e}{1}{m}$ be a base of $E_x$ and set $e_i(t) := P_t(e_i)$, then $e_1(t),\ldots,e_m(t)$ is a base of $E_{\gamma(t)}$. 
    We can write 
    \begin{equation*}
        s(\gamma(t)) = \sum_{i=1}^m s^i(t) e_i(t),
    \end{equation*}
    Then 
    \begin{equation*}
        \nabla_X s = \nabla_{\gamma'(0)} s = \sum_{i=1}^m \frac{\md s^i}{\md t}(0) e_i.
    \end{equation*}
    On the other hand, 
    \begin{equation*}
        P_t{}^{-1}(s(\gamma(t))) = \sum_{i=1}^m s^i(t) e_i,
    \end{equation*}
    so
    \begin{equation*}
        \left.\frac{\md }{\md t}\right|_{t=0} P_t{}^{-1}(s(\gamma(t))) = \sum_{i=1}^m \frac{\md s^i}{\md t}(0) e_i = \nabla_X s. \qedhere
    \end{equation*}
\end{proof}

\begin{proposition}
    Let $P^\gamma: E_{\gamma(a)} \to E_{\gamma(b)}$ be the parallel transport along $\gamma$, 
    then $P^\gamma$ satisfies the following properties:
    \begin{itemize}
        \item If the end point of $\gamma_1$ is the starting point of $\gamma_2$, 
        we can concatenate $\gamma_1$ and $\gamma_2$ to get $\gamma_1*\gamma_2$, then
        \begin{flalign*}
            &\text{(P1)} && P^{\gamma_1*\gamma_2} = P^{\gamma_2} \circ P^{\gamma_1}.& \\
        \end{flalign*} 
        \item If $\bar{\gamma}$ is the reversed curve of $\gamma$, then 
        \begin{flalign*}
            &\text{(P2)} && P^{\bar{\gamma}} = (P^\gamma)^{-1}. &
        \end{flalign*}
        \item Let $\mathcal{U}$ be an open subset in $\mathbb{R}^k$ and $\gamma: \mathcal{U}\times[0,1]\to M$ is a smooth map. 
        For any $u\in \mathcal{U}$, we denote $\gamma_u(t) := \gamma(u,t)$, then $\{\gamma_u\}$ is a family of curves smoothly depending on $u$. 
        Let $s: \mathcal{U} \to E$ be a smooth map such that $s(u)\in E_{\gamma_u(0)}$. 
        Then the map $\mathcal{U} \to E,\, u\mapsto P^{\gamma_u}(s(u))$ is also smooth. 
        For simplicity, we state this conclusion as
        \begin{flalign*}
            &\text{(P3)} && P^{\gamma} \text{ depends smoothly on } \gamma. &
        \end{flalign*}
        \item The last property is a bit technical. Let $\gamma_1,\gamma_2:[0,1]\to M$ be two curves such that $\gamma_1(0) = \gamma_2(0)$, 
        and $\gamma_1'(0) = \gamma_2'(0)$, Let $P^{\gamma_i}_t$ be the parallel transport along $\gamma_i|_{[0,t]}$, then
        \begin{equation*}
            \left.\frac{\md }{\md t}\right|_{t=0} P^{\gamma_1}_t = \left.\frac{\md }{\md t}\right|_{t=0} P^{\gamma_2}_t.
        \end{equation*}
        For simplicity, we state this conclusion as
        \begin{flalign*}
            &\text{(P4)} && \text{If } \gamma_1, \gamma_2 \text{ have contact of order } 1, \text{ then } P^{\gamma_1}, P^{\gamma_2} \text{ have contact of order } 1. &
        \end{flalign*}
    \end{itemize}
\end{proposition}
\begin{proof}
    (P1) and (P2) are straightforward from the uniqueness of solution of ODEs. 

    (P3) relies on the smooth dependence of solutions of ODEs on initial conditions and parameters.

    For (P4), given any $v\in E_{\gamma_1(0)}$, $\tilde{\gamma}_{i,v}(t):=P^{\gamma_i}_t(v)$ 
    is the horizontal lift of $\gamma_i$ starting at $v$. 
    The proposition is equivalent to saying that $\tilde{\gamma}_{1,v}'(0) = \tilde{\gamma}_{2,v}'(0)$, 
    which can be verified by calculating under a local trivialization. 
\end{proof}

\subsubsection*{Parallel Transport determines Covariant Derivative}
\begin{theorem}
    Let $E$ be a vector bundle over $M$.
    Suppose there is a family of isomorphisms 
    \begin{equation*}
        \Big\{\Phi^\gamma: E_{\gamma(0)} \to E_{\gamma(1)} \,\Big|\, \gamma:[0,1] \to M\Big\}
    \end{equation*}
    satisfying (P1)-(P4), then there exists a unique connection $\nabla$ on $E$ such that 
    $P^\gamma = \Phi^\gamma$ for any curve $\gamma$.
\end{theorem}
\begin{proof}
    For any vector field $X\in C^\infty(M,TM)$ and section $s\in C^\infty(M,E)$, at $x\in M$, 
    we can find a curve $\gamma: (-\varepsilon,\varepsilon) \to M$ 
    such that $\gamma(0) = x$ and $\gamma'(0) = X_x$.
    We define $(\nabla_{X} s)(x)$ by the formular \eqref{eq:con-def-by-para}, i.e.,
    \begin{equation*}
        (\nabla_{X} s)(x) := \left.\frac{\md }{\md t}\right|_{t=0} (\Phi^{\gamma}_t)^{-1}(s(\gamma(t))).
    \end{equation*}
    By (P4), this definition does not depend on the choice of $\gamma$. By (P3), $\nabla_X s$ is a smooth section of $E$. 

    We can verify that $\nabla$ satisfies (C2) and (C3) by direct calculation. For (C1), we need to verify that 
    $\nabla_{X_x} s$ is $\mathbb{R}$-linear in $X_x$. Let $\varphi(X_x)= \nabla_{X_x} s$, then $\varphi$ is a map from $T_x M$ to $E_x$. 
    For any $\lambda\in \mathbb{R}$, set $\gamma_\lambda(t)=\gamma(\lambda t)$, then $\gamma_\lambda(0) = x$ and $\gamma_\lambda'(0) = \lambda X_x$, so
    \begin{align*}
        \varphi(\lambda X_x) =& \nabla_{\gamma_\lambda'(0)} s = \left.\frac{\md }{\md t}\right|_{t=0} (\Phi^{\gamma_\lambda}_t)^{-1}(s(\gamma_\lambda(t))) \\
        =& \left.\frac{\md }{\md t}\right|_{t=0} (\Phi^{\gamma}_{\lambda t})^{-1}(s(\gamma(\lambda t))) \\
        =& \lambda \left.\frac{\md }{\md t}\right|_{t=0} (\Phi^{\gamma}_{t})^{-1}(s(\gamma(t))) \\
        =& \lambda\nabla_{\gamma'(0)} s = \lambda \varphi(X_x).
    \end{align*}
    It is easy to verify that $\varphi$ is additive, thus $\varphi$ is linear.

    Finally, we can verify that $P^\gamma = \Phi^\gamma$ for any curve $\gamma$ by direct calculation. 
\end{proof}

Above theorem actually says that the covariant derivative $\nabla$ is the \textcolor{red}{infinitesimal} version of the parallel transport $P^\gamma$.

\subsubsection*{Parallel Transport and Horizontal Distribution}
Let $\pi: E \to M$ be a vector bundle. 
Given a family of parallel transports $\{P^\gamma\}$, we can define a horizontal distribution $\mathcal{H}$ by
\begin{equation*}
    \mathcal{H}_p := \left\{\left.\frac{\md }{\md t}\right|_{t=0}P^{\gamma}_t(p)\,\Big|\, \gamma: [0,1] \to M \text{ with } \gamma(0) = p \right\}
    \subset T_p E.
\end{equation*}
Varying $p\in E$, we get a distribution $\mathcal{H}$ of $E$. Note that this definition is consistent with the previous one. Thus we have 
\begin{theorem}
    If $\{P^\gamma\}$ is a family of parallel transports satisfying (P1)-(P4), then $\mathcal{H}$ satisfies (H1)-(H3), 
    thus $\mathcal{H}$ is a horizontal distribution of $E$.
\end{theorem}

One can view $\mathcal{H}$ as \textcolor{red}{the collection of all directions one can possibly move via parallel transports.}

Conversely, given a horizontal distribution $\mathcal{H}$, we can define the horizontal lift, then the parallel transport. 
Expilicitly, for any curve $\gamma: [a,b] \to M$ and any point $v\in E_{\gamma(a)}$, 
there exists a unique curve $\tilde{\gamma}^{\mathcal{H}}_v: [a,b] \to E$ such that 
\begin{equation*}
    \pi(\tilde{\gamma}^{\mathcal{H}}_v(t)) = \gamma(t), \quad \tilde{\gamma}^{\mathcal{H}}_v(a) = v, \quad \tilde{\gamma}^{\mathcal{H}}_v{}'(t) \in \mathcal{H}_{\tilde{\gamma}^{\mathcal{H}}_v(t)}.
\end{equation*}
The last condition can be translated into an ODE. Since we have a decomposition $TE = \mathcal{H} \oplus \mathcal{V}$, 
we can denote the projection to $\mathcal{V}$ by $v$. Then 
\begin{equation*}
    \tilde{\gamma}^{\mathcal{H}}_v{}'(t) \in \mathcal{H}_{\tilde{\gamma}^{\mathcal{H}}_v(t)} \iff v(\tilde{\gamma}^{\mathcal{H}}_v{}'(t)) = 0,
\end{equation*}
which can be written as a linear ODEs under a local trivialization. Thus we can solve $\tilde{\gamma}^{\mathcal{H}}_v$ by the theory of ODEs. 
Now we can define $P^\gamma(v) := \tilde{\gamma}^{\mathcal{H}}_v(b)$. 
\begin{proposition}
    Fix a curve $\gamma:[0,1]\to M$ connecting $x$ and $y$, for any $v,w\in E_x$, $\lambda\in \mathbb{R}$, 
    \begin{gather*}
        \tilde{\gamma}^{\mathcal{H}}_{v+w}(t) = \tilde{\gamma}^{\mathcal{H}}_v(t) + \tilde{\gamma}^{\mathcal{H}}_w(t), \\
        \tilde{\gamma}^{\mathcal{H}}_{\lambda v}(t) = \lambda \tilde{\gamma}^{\mathcal{H}}_v(t).
    \end{gather*}
    Thus $P^\gamma: E_x \to E_{\gamma(1)}$ is a linear map. Moreover, $P^\gamma$ is an isomorphism with inverse $P^{\bar{\gamma}}$ 
    where $\bar{\gamma}$ is the reversed curve of $\gamma$.
\end{proposition}
\begin{proof}
    The first two equalities are implied by the linearity of $\mathcal{H}$, and the last one can be verified by direct calculation.
\end{proof}
\begin{theorem}
    If $\mathcal{H}$ is a horizontal distribution of $E$, then the parallel transport defined above satisfies (P1)-(P4).
\end{theorem}
\begin{proof}
    The proof is similar to the previous one.
\end{proof}
In summary, three different concepts, i.e., covariant derivative $\nabla$, parallel transport $\{P^\gamma\}$ 
and horizontal distribution $\mathcal{H}$, are equivalent to each other.
\begin{equation*}
    \text{covariant derivative} \;\Longleftrightarrow\; \text{parallel transport} \;\Longleftrightarrow\; \text{horizontal distribution}.
\end{equation*}
In fact, both $\nabla$ and $\mathcal{H}$ can be viewed as the \textcolor{red}{infinitesimal} version of $\{P^\gamma\}$.

\newpage
\subsection*{Exercises}
\addcontentsline{toc}{subsection}{Exercises}
\begin{exercise}
    Given a connection $\nabla$, we can define the covariant derivative of a $1$-form $\omega$ by
    \begin{equation*}
        (\nabla_X \omega)(Y) := X(\omega(Y)) - \omega(\nabla_X Y).
    \end{equation*}
    For an $(r,s)$-tensor field $T$, the covariant derivative $\nabla_Z T$ is 
    \begin{align*}
        (\nabla_Z T)(X_1,\ldots,X_r,\omega^1,\ldots,\omega^s):&= Z(T(X_1,\ldots,X_r,\omega^1,\ldots,\omega^s)) \\ 
        &\;\;- \sum_{i=1}^r T(X_1,\ldots,\nabla_Z X_i,\ldots,X_r,\omega^1,\ldots,\omega^s) \\
        &\;\;- \sum_{j=1}^s T(X_1,\ldots,X_r,\omega^1,\ldots,\nabla_Z \omega^j,\ldots,\omega^s).
    \end{align*}
    If we rearrange the terms, we see that this is precisely the Leibniz rule. 

    The curvature operator can also act on general tensors
    \begin{equation*}
        R(X,Y)T:= \nabla_X \nabla_Y T - \nabla_Y \nabla_X T - \nabla_{[X,Y]} T.
    \end{equation*}
    We define the covariant differential $\nabla_{\text{tensor}} T$ as the $(r,s+1)$-tensor field
    \begin{equation*}
        (\nabla_{\text{tensor}} T)(Z,\cdots,X_i,\cdots,\cdots,\omega^j,\cdots)
        := (\nabla_Z T)(\cdots,X_i,\cdots,\cdots,\omega^j,\cdots)
    \end{equation*}
    This is another way to extend the covariant derivative to tensor fields, 
    and we denote it by $\nabla_{\text{tensor}}$ to distinguish it from the exterior covariant derivative.

    \textcolor{blue}{Prove that} for a torsion-free connection, $\nabla_{\text{tensor}}$ satisfies \textbf{the Ricci identity}
    \begin{equation*}
        (\nabla_{\text{tensor}}^2 T)(X,Y,\cdots,\cdots) - (\nabla_{\text{tensor}}^2 T)(Y,X,\cdots,\cdots) = (R(X,Y)T)(\cdots,\cdots)
    \end{equation*}
\end{exercise}

% \begin{proof}
%     For an $(r,s)$-tensor field $T$, we have
%     \begin{align*}
%         &(\nabla_{\text{tensor}}^2 T)(X,Y,\cdots,\cdots) = (\nabla_X (\nabla_{\text{tensor}} T))(Y,\cdots,\cdots) \\
%         =& X((\nabla_{\text{tensor}} T)(Y,\cdots,\cdots)) - (\nabla_{\text{tensor}} T)(\nabla_X Y,\cdots,\cdots) \\
%         &- \sum_{i=1}^r (\nabla_{\text{tensor}} T)(Y,\cdots,\nabla_X X_i,\cdots,\cdots) 
%         - \sum_{j=1}^s (\nabla_{\text{tensor}} T)(Y,\cdots,\cdots,\nabla_X \omega^j,\cdots) \\
%         =& X((\nabla_Y T)(\cdots,\cdots)) - (\nabla_{\nabla_X Y} T)(,\cdots,\cdots) \\
%         &- \sum_{i=1}^r (\nabla_Y T)(\cdots,\nabla_X X_i,\cdots,\cdots) 
%         - \sum_{j=1}^s (\nabla_Y T)(\cdots,\cdots,\nabla_X \omega^j,\cdots) \\
%         =& (\nabla_X \nabla_Y T)(\cdots,\cdots) - (\nabla_{\nabla_X Y} T)(,\cdots,\cdots),
%     \end{align*}
%     Thus 
%     \begin{align*}
%         &(\nabla_{\text{tensor}}^2 T)(X,Y,\cdots,\cdots) - (\nabla_{\text{tensor}}^2 T)(Y,X,\cdots,\cdots) \\
%         =& (\nabla_X \nabla_Y T)(\cdots,\cdots) - (\nabla_Y \nabla_X T)(\cdots,\cdots) - (\nabla_{[X,Y]} T)(\cdots,\cdots) \\
%         =& (R(X,Y)T)(\cdots,\cdots).
%     \end{align*}
% \end{proof}

\begin{exercise}
    For a vector field $Z\in C^\infty(M,TM)$, suppose $\nabla$ is torsion-free. 
    Both $\nabla^2 Z$ and $\nabla_{\text{tensor}}^2 Z$ are $(1,2)$-tensor fields, 
    but by Ricci identity, we have
    \begin{equation*}
        (\nabla_{\text{tensor}}^2 Z)(X,Y) - (\nabla_{\text{tensor}}^2 Z)(Y,X) = R(X,Y)Z = (\nabla^2 Z)(X,Y).
    \end{equation*}
    Can you explain this phenomenon?
\end{exercise}

\begin{exercise}[Torsion 2-form]
    Let $\nabla$ be a connection on $TM$, $\{e_i\}$ be a local frame of $TM$, 
    $\{\theta^i\}$ be the corresponding dual coframe 
    and $\omega$ be the connection 1-form of $\nabla$ w.r.t. $\{e_i\}$. 
    Define 
    \begin{equation*}
        \tau^i := \md\theta^i + \omega_j^i \wedge \theta^j.    
    \end{equation*}
    If we arrange $\theta=(\ftu{\theta}{1}{n})^T$ as a column vector, then the above formula can be written as
    \begin{equation*}
        \tau = \md\theta + \omega \wedge \theta.
    \end{equation*}
    then $\tau$ is called the \textbf{torsion 2-form} of $\nabla$ w.r.t. $\{e_i\}$. 

    Let $T(e_i,e_j) = T_{ij}^k e_k$, where $T$ is the torsion of $\nabla$. \textcolor{blue}{Prove that} 
    \begin{equation*}
        \tau^k(e_i,e_j) = T_{ij}^k.
    \end{equation*}
    Thus 
    \begin{equation*}
        \tau^k = \frac{1}{2} T_{ij}^k \theta^i \wedge \theta^j.
    \end{equation*}
\end{exercise}

% \begin{proof}
%     Since $\tau^k = \md\theta^k + \omega_l^k \wedge \theta^l$, 
%     \begin{align*}
%         &\tau^k(e_i,e_j) e_k = (\md\theta^k)(e_i,e_j)e_k + (\omega_l^k \wedge \theta^l)(e_i,e_j)e_k \\
%         =& e_i(\theta^k(e_j))e_k - e_j(\theta^k(e_i))e_k - \theta^k([e_i,e_j])e_k + \omega_l^k(e_i)\theta^l(e_j)e_k - \omega_l^k(e_j)\theta^l(e_i)e_k \\
%         =& -\theta^k([e_i,e_j])e_k + \omega_j^k(e_i)e_k - \omega_i^k(e_j)e_k \\
%         =& \nabla_{e_i} e_j - \nabla_{e_j} e_i-[e_i,e_j]  = T(e_i,e_j).
%     \end{align*}
% \end{proof}

\begin{exercise}[Levi-Civita connection]
    Let $(M,g)$ be a Riemannian manifold, $\nabla$ be a connection on $TM$. Show that 
    \begin{enumerate}
        \item $\nabla$ is torsion-free if and only if $\tau=0$, i.e., 
        \begin{equation*}
            \md\theta^i + \omega_j^i \wedge \theta^j = 0.
        \end{equation*}
        \item If $\{e_i\}$ is a local orthonormal frame, 
        then $\nabla$ is compatible with $g$ if and only if $\omega$ is skew-symmetric, i.e.,
        \begin{equation*}
            \omega_j^i + \omega_i^j = 0.
        \end{equation*}
        \item Under a local orthonormal frame, there exists a unique connection 1-form $\omega$ satisfying 
        \begin{equation*}
            \md\theta^i + \omega_j^i \wedge \theta^j = 0, \quad
            \omega_j^i + \omega_i^j = 0.
        \end{equation*}
        This is the existence and uniqueness of the Levi-Civita connection.
    \end{enumerate} 
\end{exercise}

\begin{exercise}
    Let $\nabla$ be a connection on $TM$, $\{e_i\}$ be a local frame of $TM$, 
    $\{\theta^i\}$ be the corresponding dual coframe 
    and $\omega$ be the connection 1-form of $\nabla$ w.r.t. $\{e_i\}$. 
    The following two equations 
    \begin{align*}
        \md\theta =& \tau - \omega \wedge \theta, \\
        \md\omega =& \Omega - \omega \wedge \omega ,
    \end{align*}
    are called the \textbf{structure equations}. Show that $\tau$ and $\Omega$ satisfy
    \begin{align*}
        \md\tau =& \Omega \wedge \theta - \omega \wedge \tau, \\
        \md\Omega =& \Omega \wedge \omega - \omega \wedge \Omega.
    \end{align*}
    These two equations are called the \textbf{Bianchi identities}. Furthermore, show that 
    \begin{equation*}
        \md(\Omega^k) = \Omega^k \wedge \omega - \omega \wedge \Omega^k.
    \end{equation*}
\end{exercise}

\begin{exercise}
    Show that if $\nabla$ is torsion-free, then $\md\tau=\Omega\wedge\theta-\omega\wedge\tau$ 
    is equivalent to \textbf{the first Bianchi identity}
    \begin{equation*}
        R(X,Y)Z + R(Y,Z)X + R(Z,X)Y = 0.
    \end{equation*}
\end{exercise}

\begin{exercise}
    Show that if $\nabla$ is torsion-free, then $\md\Omega=\Omega\wedge\omega-\omega\wedge\Omega$ 
    is equivalent to \textbf{the second Bianchi identity}
    \begin{equation*}
        (\nabla_X R)(Y,Z) + (\nabla_Y R)(Z,X) + (\nabla_Z R)(X,Y) = 0.
    \end{equation*}
\end{exercise}

\begin{exercise}
    For $\{e_i\}=\{\frac{\partial}{\partial x^i}\},\,\{\tilde{e}_j\}=\{\frac{\partial}{\partial \tilde{x}^j}\},
    \,\omega^k_j=\Gamma^{k}_{ij}\md x^i$, show that the transformation formula 
    \begin{equation*}
        \tilde{\omega} = a^{-1} \md a + a^{-1} \omega a
    \end{equation*}
    reduces to the well-known transformation formula of Christoffel symbols:
    \begin{equation*}
        \tilde{\Gamma}_{ij}^k 
        = \pd{\tilde{x}^k}{x^l}\pd{x^m}{\tilde{x}^i}\pd{x^n}{\tilde{x}^j}\Gamma_{mn}^l 
        + \frac{\partial^2 x^l}{\partial \tilde{x}^i \partial \tilde{x}^j} \frac{\partial \tilde{x}^k}{\partial x^l}.
    \end{equation*}
\end{exercise}

\begin{exercise}
    Let $\nabla$ be a connection on $TM$. It determines a horizontal distribution $\mathcal{H}$ on $TM$. 
    Then $T(TM) = \mathcal{V}\oplus\mathcal{H}$, where $\mathcal{V}$ is the vertical bundle of $TM$. 
    Show that $\nabla$ is torsion-free if and only if for any $X,Y\in C^\infty(M,TM)$,
    \begin{equation*}
        [X^h,Y^h] = [X,Y]^h.
    \end{equation*}
\end{exercise} 